% ============================================================
% DRAFT NOTE (discussion)
% Contains: Interpretation by compressor and ratio, limits of the evidence,
% and a short conclusion.
% Written now: Scope limitations and the conditional interpretation frame.
% Pending: The result-grounded account of success or failure by compressor and
% ratio.
% Sources: docs/goal.md; docs/methodology.md; docs/related_work.md
% ============================================================
\section{Discussion}

The final discussion will interpret forecasting performance separately for
each compressor and removal ratio, with the preregistered decision gate as the
controlling evidence. \todo{Describe which compressors and ratios support or
do not support magnitude forecasting, based on the held-out results.}

\subsection{Limitations}

This initial study concerns Llama-3.1-8B-Instruct \citep{llama31} on the
single task IFEval. Its claim is limited to known compressors at known
supported removal ratios. It does not establish that the features transfer to
an unseen compressor, another model, another task, or an unseen ratio. IFEval
loose accuracy is task grounded and deterministic, but it is still a bounded
measure of instruction following rather than universal response quality.

The parity ladder constrains the interpretation of every label: live-fork and
matched-reprefill targets are comparable only to the extent established by the
reported parity evidence. \todo{Add any result-specific limitations revealed
by calibration, positive-damage slices, or ratio-level outcomes.}

\subsection{Conclusion}

HERALD frames KV-cache compression damage as a paired final-quality effect at
a pre-switch decision boundary and tests whether causal state statistics can
forecast its signed magnitude. The resulting evidence will state the claim at
the compressor and ratio level, rather than relying on pooled accuracy.
Future work may examine how a validated forecast can inform downstream
serving policies.
